\section{Truth-table companion syntax}

Truth tables are metadata paired to a process; they are not \code{.qpucir}
commands. A protocol describes \emph{how} a circuit works; its companion
\code{.qpuio} describes selected input/output expectations.

\subsection{Text \code{.qpuio} structure}

A text table may use spaces or commas:

\begin{lstlisting}
MAIN-PROCESS: HalfAdder
INPUTS: A B
OUTPUTS: Carry Sum
# A B Carry Sum
0 0p 0p 0p 0p
1 0p 1p 0p 1p
2 1p 0p 0p 1p
3 1p 1p 1p 0p
\end{lstlisting}

\subsection{\code{MAIN-PROCESS:} table header}

The colon is required in the table format:
\begin{lstlisting}
MAIN-PROCESS: HalfAdder
\end{lstlisting}
This name associates metadata with its protocol. Notice the deliberate syntax
difference: protocol source writes \code{MAIN-PROCESS HalfAdder}, while the
truth-table header writes \code{MAIN-PROCESS: HalfAdder}.

\subsection{\code{INPUTS:} and \code{OUTPUTS:} column groups}

\begin{lstlisting}
INPUTS: A B
OUTPUTS: Carry Sum
\end{lstlisting}
These lines divide the columns into values supplied before simulation and
values checked after simulation. Their order must match the later column header.
INPUTS should correspond to state-typed PARAMS; OUTPUTS should correspond to
RETURNVALS. OUTPUTS must contain at least one name.

Explicit groups remove ambiguity. If they are absent, the loader can derive the
split only when a matching protocol is available and the combined PARAMS and
RETURNVALS widths match the table header.

\subsection{Column header row}

\begin{lstlisting}
# A B Carry Sum
\end{lstlisting}
The first nonmetadata row beginning with a hash and whitespace or comma is the
ordered column header. Here A and B are inputs, followed by Carry and Sum
outputs. Unlike an ordinary protocol hash comment, this row has structural
meaning inside a \code{.qpuio} file.

\subsection{Indexed data rows}

\begin{lstlisting}
0 0p 0p 0p 0p
1 0p 1p 0p 1p
\end{lstlisting}
The first cell is a decimal row index. Indices begin at zero and increase by
one without gaps. Remaining cells align left-to-right with the column header.
Every row must have exactly the expected width, and every state cell must be
\code{0p}, \code{1p}, or \code{sp}.

For $n$ binary input columns, a complete table contains $2^n$ rows. The usual
ordering interprets input columns most-significant first: for two inputs the
sequence is 00, 01, 10, 11. The loader checks sequential row indices, while
semantic validation also rejects duplicate input patterns.

\subsection{CSV spelling}

Commas may replace whitespace consistently:
\begin{lstlisting}
MAIN-PROCESS: HalfAdder
INPUTS: A,B
OUTPUTS: Carry,Sum
#,A,B,Carry,Sum
0,0p,0p,0p,0p
\end{lstlisting}
Spaces around commas are trimmed. Do not mix extra prose into a data row;
trailing hash comments are allowed only where the table loader can strip them
after the data content.

\subsection{Partial truth tables}

A partial table lists selected cases rather than all $2^n$ combinations. It
must contain at least one row, no duplicate input pattern, and no more than
$2^n$ rows. Partial tables are useful for state-machine transitions or focused
regressions, but passing one proves behavior only for the listed cases.

\begin{center}
\begin{tabular}{c|c|c}\toprule
Input count $n$ & Complete row count $2^n$ & Example\\\midrule
0&1&constant-output process\\
1&2&0, 1\\
2&4&00, 01, 10, 11\\
3&8&000 through 111\\
4&16&0000 through 1111\\\bottomrule
\end{tabular}
\end{center}

\subsection{JSON \code{.qpuio} envelope}

The equivalent machine-oriented JSON form is:
\begin{lstlisting}[language={}]
{
  "format": "qpuio",
  "version": 1,
  "processName": "HalfAdder",
  "inputColumns": ["A", "B"],
  "outputColumns": ["Carry", "Sum"],
  "rows": [
    ["0p", "0p", "0p", "0p"],
    ["0p", "1p", "0p", "1p"],
    ["1p", "0p", "0p", "1p"],
    ["1p", "1p", "1p", "0p"]
  ]
}
\end{lstlisting}

The \code{format} value must be \code{qpuio} and \code{version} must be 1.
Column names must be nonempty strings, outputColumns must not be empty, and
every nested row array must have exactly
\[
\text{inputColumns.length}+\text{outputColumns.length}
\]
valid cells.

\subsection{How a table test runs}

For each listed row, the runner:
\begin{enumerate}
  \item maps input cells to state PARAMS by column name and order;
  \item creates the circuit's start-state vector;
  \item executes every compiled gate, including child expansions;
  \item measures the declared RETURNVALS;
  \item compares only those measured outputs with expected output cells.
\end{enumerate}
An expected output \code{sp} accepts either measured bit. An input \code{sp}
initializes an equal plus-state superposition. Because one probabilistic sample
cannot establish a full distribution, tables are strongest for deterministic
basis-state behavior.

\subsection{Protected bundled tables}

The bundled SingleBitFullAdder, TwoBitFullAdder, FourBitFullAdder, and PhaseDemo
tables are canonical site assets. Do not rewrite them to make an experiment
pass. Create a new uploaded or corrected process/table pair when testing
different expectations. This keeps examples reproducible and prevents a
protocol correction from silently changing its specification.

